En una càpsula de Petri plena d'aigua destil·lada hem submergit dues plaques metàl·liques paral·leles connectades a una diferència de potencial de 12,0~V, tal com mostra la figura. Les dues plaques estan separades per una distància de 6,00~cm. Amb un voltímetre, explorem la diferència de potencial entre la placa negativa i diferents punts de la regió intermèdia.

\figura[0.35]{fig1.png}

\begin{apartat}{1,25}
Calculeu el camp elèctric (suposant que és uniforme) entre les dues plaques, i indiqueu-ne també la direcció i el sentit. Feu un dibuix en què representeu, de manera aproximada, les superfícies equipotencials que espereu trobar a la regió compresa entre les dues plaques i indiqueu el valor del potencial en cadascuna de les superfícies representades.
\end{apartat}

\begin{apartat}{1,25}
Amb la sonda, tal com veiem a la figura, el voltímetre indica 7,0~V. Calculeu el treball que hauria de fer una força externa per a desplaçar una càrrega positiva de 0,1~$\mu$C des d'aquest punt fins a la placa positiva.
\end{apartat}
